
The semiconductor industry keeps seeking for reducing the design time and efforts in modern integrated circuit implementation, 
while detailed routing is a dead-or-alive critical element in design automation tooling for advanced node enablement, 
which is known as the most time-consuming step in the back-end design flow due to its high complexity and enormous solution space for the very large-scale integration (VLSI) routing problem.

In this paper, we propose a FPGA-based implementation to accelerate the detailed routers, uses efficient memory access management and parallelism to achieve better performance.
Experimental results demonstrate that detailed routers obtain stable and effective acceleration on FPGA. 
Our design gets 2.1$\times$-2.3$\times$, 2.1$\times$-2.5$\times$, 1.3$\times$ speedup compared to TritonRoute, Dr.CU,and current optimum routing kernel acceleration work, 
also obtains 20\% speedup for  detailed routing flow, and 1.0$\times$-1.2$\times$ resource utilization compared to TritonRoute.
% \textbf{(1) Incompatibility between element-wise operations and Tensor Core.} 
% Linear operations (matrix multiplications) and element-wise operations are the two dominating operations in Mamba.
% The time proportion of element-wise operations escalates significantly (\textit{e.g.}, $>$60\% with 2048 input length). 
% These operations do not need reduction, which is not compatible with the existing Tensor Core-based architectures (\textit{e.g.}, 1/16 normalized speed). 
% \textbf{(2) Large area overhead for nonlinear function unit.} The optimized nonlinear function unit like exponential unit still occupies $>$30\% of the processing element (PE) area.
% % , which introduces large area and power overheads.
% \textbf{(3) Large memory access but limited data sharing for element-wise operations.} 
% % Linear and element-wise operations also show different characteristics of data sharing.
% Linear and element-wise operations in Mamba exhibit large compute intensity variance (\textit{e.g.}, $\sim$3 orders of magnitude) and large read/write ratio variance (\textit{e.g.}, $>$3 orders). 
% Due to the limited data sharing in element-wise operations, it is useless to apply the existed methods like tiling to element-wise operations.

% In response to these challenges, we propose a Mamba accelerator with reconfigurable architecture, \we.
% % , controlled by instruction with reconfigurable computing units (CUs) and processing elements (PEs).
% Then, we propose three novel approaches in this paper.
% \textbf{(1) Reduction alternative PE array architecture} for both linear and element-wise operations. 
% For linear operations, the reduction tree connected to PE arrays is enabled and executes the reduction operation.
% For element-wise operations, the reduction tree is disabled and the output bypasses.
% \textbf{(2) Reusable nonlinear function unit} based on the reconfigurable PE. 
% We decompose the exponential function into element-wise operations and a shift operation by a fast biased exponential algorithm,
% and the activation function (SiLU) into a range detection and element-wise operations by a piecewise approximation algorithm.
% Thus, the reconfigurable PEs are reused to execute nonlinear functions with negligible accuracy loss.
% \textbf{(3) Intra-operation and inter-operation buffer management strategy}.
% We propose intra-operation buffer management strategy to maximize input data sharing for linear operations within operations, and inter-operation strategy for element-wise operations between operations.
% % We conduct extensive experiments on Mamba model families with different model sizes.
% We conduct extensive experiments on Mamba model families with different sizes.
% \we achieves up to 463.22$\times$/11.66$\times$ speedup and up to 9761.42$\times$/242.52$\times$ energy efficiency compared to Intel Xeon 8358P CPU and NVIDIA Tesla A100 GPU implementations, respectively. 

